% Portiegrootte
% Not a real \chapter (it's a hand-styled heading like Register's), so it
% needs its own Inhoud/PDF-bookmark entry, same as imakeidx's intoc=true
% does for Register.
\phantomsection\addcontentsline{toc}{chapter}{Portiegrootte}
\vspace*{0.4cm}
{\handfont\fontsize{60}{60}\selectfont\color{terracotta}Portiegrootte\par}
\vspace{8pt}{\color{terracotta}\rule{2.6cm}{0.5pt}\;\raisebox{-0.2ex}{$\scriptstyle\blacklozenge$}}\par
\vspace{1.4em}

\lettrine{H}{oeveel} gram pasta reken je nu eigenlijk per persoon? Dit
boek gaat overal van dezelfde aannames uit, zodat een recept ook
duidelijk is op de punten waar de tekst zelf niets over zegt.

\blockrule{Aantal porties}
Tenzij er bij de kooktijd van een recept iets anders staat (zoals ``Voor 4
personen''), is dit boek overal doorgerekend op 5 personen. Wil je een
recept voor een ander aantal maken, schaal dan de hoeveelheden naar rato.

\blockrule{Hoeveelheden per persoon}
Dit zijn de hoeveelheden die we zelf aanhouden voor het standaard aantal
van 5 personen, hier uitgesplitst naar wat dat per persoon betekent. Handig
als je een recept wilt op- of afschalen.
\begin{ingredients}
\ing{70 g}{gedroogde pasta}
\ing{100 g}{verse pasta}
\ing{35 g}{rijst, droog gewicht, als bijgerecht}
\ing{35--65 g}{risottorijst, droog gewicht, meer als het een hoofdgerecht is}
\ing{32 g}{couscous, droog gewicht, als bijgerecht}
\ing{200 g}{aardappelen, geschild, voor een pure stamppot}
\ing{100 g}{aardappelen, geschild, in een stamppot met net zoveel andere groente}
\ing{120 g}{aardappelen, voor friet}
\end{ingredients}

\vspace{0.6em}
{\footnotesize\itshape\color{inkfaint}Dit zijn de vuistregels die we zelf
aanhouden, geen exacte wetenschap. Bij een stevige eter, of als je toch
extra voor de lunch de volgende dag wilt overhouden, reken je gerust
ruimer.\par}
